%------------------------------------------%
% Cannabis Data Science #58
% Date: 3/23/2022
%------------------------------------------%
\documentclass[xcolor={dvipsnames}]{beamer}
\hypersetup{pdfpagemode = FullScreen}
\mode<presentation>{
  \usetheme{Boadilla}
  \usecolortheme{orchid}
  \usefonttheme{default}
  \setbeamertemplate{navigation symbols}{}
  \setbeamertemplate{caption}[numbered]
}
\setbeamersize{
  text margin left = 0.5in,
  text margin right = 0.5in
}

%------------------------------------------%
% Title
%------------------------------------------%
\title[\textbf{Cannabis Data Science \#58}]{}
\author{Cannabis Data Science}
\institute[]{\Large Cannabis Data Science \#58}
\date{March \nth{23}, 2022}

%------------------------------------------%
% Packages
%------------------------------------------%
\usepackage[english]{babel}
\usepackage[utf8x]{inputenc}
\usepackage{tikz} % For styling.
\usepackage{xparse}

%------------------------------------------%
% Colors
%------------------------------------------%
\definecolor{Green}{RGB}{34, 153, 84}
\definecolor{LightGreen}{RGB}{218, 247, 166}
\definecolor{DarkGreen}{RGB}{2, 48, 32}
\definecolor{Orange}{RGB}{255, 87, 51}
\definecolor{DarkOrange}{RGB}{199, 0, 57}
\definecolor{Yellow}{RGB}{255, 195, 0}

%------------------------------------------%
% Theme
%------------------------------------------%
\setbeamercolor*{palette primary}{bg=LightGreen, fg=DarkGreen}
\setbeamercolor*{palette secondary}{bg=LightGreen, fg=DarkGreen}
\setbeamercolor*{palette tertiary}{bg=LightGreen, fg=DarkGreen}

%------------------------------------------%
% Packages
%------------------------------------------%
\usepackage{amsmath}
\renewcommand*\footnoterule{} % No separating line on footnote.
\usepackage{mathtools} % For annotating equations.
\usepackage{hhline} % for double bars.
\usepackage[super]{nth} % For formatting 1st, 2nd, 3rd, etc.
\usepackage{graphicx, caption, subcaption}

%------------------------------------------%
% Commands
%------------------------------------------%

% Top space.
\newcommand\T{\rule{0pt}{2.5ex}}

% Bottom space.
\newcommand\B{\rule[-1.25ex]{0pt}{0pt}}

% Blocks.
\newenvironment<>{Block}[2][.9\textwidth]
  {\setlength{\textwidth}{#1}
  \begin{actionenv}#3
    \def\insertblocktitle{#2}\par
    \usebeamertemplate{block begin}}
  {\par\usebeamertemplate{block end}
  \end{actionenv}}

% Balls.
\defbeamertemplate{enumerate item}{largeball}
{\begin{pgfpicture}{-1ex}{-0.65ex}{1.5ex}{1.5ex}
\usebeamercolor[fg]{item projected}
{\pgftransformscale{2.5}\pgftext{\Large\pgfuseshading{bigsphere}}}
{\pgftransformshift{\pgfpoint{0pt}{0.5pt}}
\pgftext{\usebeamerfont*{item projected}\small\insertenumlabel}}
\end{pgfpicture}}

% Fancy arrows.
\NewDocumentCommand\UpArrow{O{2.0ex} O{black}}{%
   \mathrel{\tikz[baseline] \draw [->, line width=0.5pt, #2] (0,0) -- ++(0,#1);}} % Fancy up-arrow.
\NewDocumentCommand\DownArrow{O{2.0ex} O{black}}{%
   \mathrel{\tikz[baseline] \draw [<-, line width=0.5pt, #2] (0,0) -- ++(0,#1);}} % Fancy down-arrow.

% Equations with numbers on the left.
\makeatletter
\newcommand{\LeftEqNo}{\let\veqno\@@leqno}
\makeatother

%------------------------------------------%
% Presentation
%------------------------------------------%
\begin{document}

% Title page.
\begin{frame}{}
  \includegraphics[scale=0.33]{images/logo.pdf}
  \vspace*{-2\baselineskip}
  \titlepage

  % TODO: Add flare to title page?
  % Background
%\tikz[remember picture, overlay]
%\node[opacity=1.0, inner sep=0pt] at (current page.center){
%  \includegraphics[width=\paperwidth, height=\paperheight]{images/cover.pdf}
%};  
  
\end{frame}

% Run the gamut with beverage brand analysis.

% Using beverage product type for simplicity. The analysis can be readily applied to all product types.



%------------------------------------------%
% Data Science
%------------------------------------------%
\section{Data Science}

\begin{frame}{When did data science begin?}

\vspace*{0.125\baselineskip}

% TODO: Write synopsis:
John Graunt
%https://en.wikipedia.org/wiki/John_Graunt


% Andrew S. C. Ehrenberg

% negative binomial distribution to the numbers of purchases of a brand of consumer goods.


% NBD-Dirichlet model of brand choice successfully modelled the repeated category and brand purchases within a wide variety of markets.


%Estimate {\bfseries penetration rate} (percent).
%The percentage of the relevant population that has purchased a given brand (or category) at least once in the time period under study.


\end{frame}


%------------------------------------------%
% Application to the Cannabis Industry
%------------------------------------------%
\section{Application to the Cannabis Industry}

\begin{frame}{}


% cannabis--infused beverage brands?

% pre-roll brands?

% solid edible brands? Overlap with beverage brands?


\end{frame}


%------------------------------------------%
% Takeaway
%------------------------------------------%
\section{Takeaway}
\begin{frame}{}

\begin{center}
\begin{minipage}{3.85in}

% Thank you.
\includegraphics[width=.25in]{images/prayer.png} {\Large \textbf{Thank you for coming.}}\\

\begin{center}
\begin{minipage}{.9\linewidth}
\begin{Block}{Lesson of the Day}

\vspace{0.5\baselineskip}

\begin{itemize}

\item ?

\end{itemize}

\end{Block}
\end{minipage}
\end{center}

\vfill

\end{minipage}
\end{center}

\end{frame}

%------------------------------------------%
% Fin.
%------------------------------------------%
\end{document}
